\section{Exercise 2}
\textit{let $T$ be the triangle in $\mathbb{R}^2$ with vertices (1,0), (0,1) and (-1,0) traversed 
counterclockwise. Also let $F$ be the vector field:$F(x,y) = (e^y, sin(\pi x))$ and $S$ the line integral:}
\[S = \oint_T F \cdot dr\]
\begin{itemize}
    \item Calculate the line integral S directly
    \item Calculate the line integral S using Greens Theorem
\end{itemize}

To evaluate this integral, the bound of integration must fist be defined.
This since is a triangle, in the counterclockwise direction, three line segments can be defined, 
from the given three points:
\[P_1=(1,0),\qquad P_2=(0,1),\qquad P_3=(-1,0),\qquad
C_1 = P_2 - P_1, \qquad C_2 = P_3-P_2, \qquad C_3 = P_1-P_3\]
This can written as a direction vector: $r(t) = P+t(Q-P)$ where P is the origin 
point and Q is the ending point. Then the vector equations can be split up into
the x and y components. Also note that t is bounded by $0\le t \le 1$:
%vect 1
\[
r_{C_1}(t) = P_1 +t(P_2 - P_1) \Rightarrow 
\begin{pmatrix} 
    1\\
    0 
\end{pmatrix}
+t
\begin{pmatrix}
    0-1\\
    1-0
\end{pmatrix} 
= 
\begin{pmatrix} 
    1\\
    0 
\end{pmatrix}+
t
\begin{pmatrix}
    -1\\
    1
\end{pmatrix},
\qquad r'_{C_1}(t) =
\begin{pmatrix}
    -1\\
    1
\end{pmatrix}
\]
%vect 2
\[
r_{C_2}(t) = P_2 +t(P_3 - P_2) \Rightarrow 
\begin{pmatrix} 
    0\\
    1 
\end{pmatrix}
+t
\begin{pmatrix}
    -1-0\\
    0-1
\end{pmatrix} 
= 
\begin{pmatrix} 
    0\\
    1 
\end{pmatrix}+
t
\begin{pmatrix}
    -1\\
    -1
\end{pmatrix},
\qquad r'_{C_2}(t) = 
\begin{pmatrix}
    -1\\
    -1
\end{pmatrix}
\]
%vect 3
\[
r_{C_3}(t) = P_3 +t(P_1 - P_3) \Rightarrow 
\begin{pmatrix} 
    -1\\
    0 
\end{pmatrix}
+t
\begin{pmatrix}
    1-(-1)\\
    0-0
\end{pmatrix} 
= 
\begin{pmatrix} 
    -1\\
    0 
\end{pmatrix}+
t
\begin{pmatrix}
    2\\
    0
\end{pmatrix},
\qquad r'_{C_3}(t) =
\begin{pmatrix}
    2\\
    0
\end{pmatrix}
\]
The integral then becomes the sum of the each of the integrals:
\[\oint_C \mathbf{F} \cdot dr = \sum_{i=1}^{3} \int_{0}^{1} F(r_{C_i}(t))\cdot r'_{C_i}(t)\]